% Notes-only candidate; no canonical registration before independent audits.
\subsection{An exact ghost in the full kernel of descent}

\begin{statement}{C.44}{Open}{Candidate: descent has a nonzero parity kernel}
There exists an actual multiplier-one full gentle element $e$ with
\[
             D(e)=1,\qquad \Delta(e)=1.                 \tag{C.44.1}
\]
Consequently $D$ is not injective, and the restriction of $\Delta$
to its closed normal kernel is onto $C_2$. One and the same element
satisfies
\[
       (\Delta(e),\Delta(D(e)),\Delta(D^2(e)),\ldots)
                         =(1,0,0,\ldots).              \tag{C.44.2}
\]
The induced endomorphism on the abelianization of the quotient by
the closed normal Galois image is also not injective. This does not
assert surjectivity of $D$, independence of all its iterated parity
characters, or any new statement about the pentagon subgroup.
\end{statement}

\begin{proof}[Candidate proof, pending full independent audits]
The new point is that the relative lifting argument works for a
closed normal subgroup which need not be open. We first establish
that point with all its profinite comparisons retained.

\paragraph{Relative lifting for a closed normal subgroup.}
Let $\Gamma=\widehat{C_2*C_3}=\langle s,r\mid s^2=r^3=1\rangle$,
let $\tau=sr$, and let $U$ be any closed normal subgroup of $\Gamma$.
Write $R_{\rm nil}(U)$ for the intersection of the kernels of its
finite nilpotent quotients, and $O^p(U)$ for the kernel of its
maximal pro-$p$ quotient. Both are characteristic in $U$, hence
normal in $\Gamma$.

For any profinite group $U$, with closed lower-central subgroups,
\[
 R_{\rm nil}(U)=\bigcap_{n\ge1}\gamma_n(U),\qquad
 q(R_{\rm nil}(U))=\gamma_\infty(q(U))                  \tag{C.44.3}
\]
for every finite quotient map $q$ of $\Gamma$. For the first
equality, each bounded-class quotient is an inverse limit of finite
nilpotent groups, and every finite nilpotent group has bounded
class. For the second, $q(\gamma_n(U))=\gamma_n(q(U))$ by
continuity, density of commutator generators and finite image.
Images commute with this descending intersection: a compatible
nonempty sequence of fibres inside the compact groups
$\gamma_n(U)$ has nonempty intersection. The finite lower-central
sequence stabilizes. No finite generation or openness of $U$ is used.

Every peripheral automorphism of
$\Gamma/R_{\rm nil}(U)$ with three specified positive conjugators
lifts to a full peripheral automorphism of $\Gamma$, retaining
those conjugators. Indeed, take any finite source
$q_E:\Gamma\twoheadrightarrow E$ and put
$A=\gamma_\infty(q_E(U))$. Then $A$ is normal and
$[q_E(U),A]=A$. By (C.44.3), $E/A$ is an actual quotient of
$\Gamma/R_{\rm nil}(U)$. In the finite scalar-cover construction
of P.3, every auxiliary source homomorphism to a central extension
of $E/A$ has nilpotent image of $U$: it is a subgroup of a central
extension of the finite nilpotent group $q_E(U)/A$. A central
extension increases nilpotency class by at most one. Thus every
auxiliary source homomorphism kills $R_{\rm nil}(U)$.

Precomposing these actual source homomorphisms with the prescribed
automorphism supplies the central comparison required by P.3.
Its trace count, subgroup-by-subgroup comparison and M\"obius
inversion require only normality of the finite $A$. They therefore
give a generating map into each $E$ with the required retained
frame values. Any finite collection of conditions on the retained
profinite quotient and its three conjugators is contained in one
finite joint source $E$. Compactness in $\Gamma^3$ gives all of
them simultaneously. Generation of every finite source quotient
makes the resulting endomorphism onto, hence an automorphism of
the finitely generated profinite group $\Gamma$. This repeats the
complete finite refinement of P.3; $\Gamma/U$ was nowhere required
to be finite. A frame prescribed to be one can finally be made
literally one by an inner adjustment in $R_{\rm nil}(U)$.

The exact relative fibre product used below also holds for this
closed $U$:
\[
 \Gamma/R_{\rm nil}(U)
       \simeq\prod_{p,\,\Gamma/U}\Gamma/O^p(U).          \tag{C.44.4}
\]
On the normal kernel this is
$U/R_{\rm nil}(U)\simeq\prod_p U^{(p)}$. To see its
surjectivity, every finite collection of finite $p$-group
quotients has product image: a subdirect product of groups of
pairwise coprime orders is their whole product. Compactness gives
the full product. Its kernel is exactly the intersection of all
finite nilpotent quotient kernels. For (C.44.4), choose a lift in
$\Gamma$ of the common $\Gamma/U$ coordinate and reduce the
remaining tuple to this kernel statement. This proves both
surjectivity and the asserted topological isomorphism.

\paragraph{A word killed by all six full fillings.}
Put $F=\ker(\Gamma\to S_3)$ and use the temporary marking of C.23,
\[
 x=\tau^{-2},\qquad y=sxs,\qquad z=(xy)^{-1}.
\]
The actual filling of the question is, in these letters,
\[
 \varphi(x)=\tau^{-1},\qquad\varphi(y)=s,\qquad
                         \varphi(z)=r.                \tag{C.44.5}
\]
Let $\rho=\operatorname{Int}(r)$ and
$\theta=\operatorname{Int}(s)$ on $F$. They satisfy
$\rho(x,y,z)=(y,z,x)$ and $\theta(x,y)=(y,x)$.
Take the six actual representatives
$1,\rho,\rho^2,\theta,\rho\theta,\rho^2\theta$ and set
\[
 H=F'F^3,\qquad
 U=H\cap\bigcap_\sigma\ker(\varphi\sigma:F\to\Gamma).
                                                               \tag{C.44.6}
\]
This is a closed normal subgroup of $\Gamma$. Indeed $H$ is
characteristic in $F$, and composition of the chosen representatives
differs from another by an inner automorphism of $F$. Normality
of each kernel in $F$ makes their intersection invariant under
the actual $\Gamma$ action. No assertion that $U$ is open is made.

Use the commutator convention $[a,b]=aba^{-1}b^{-1}$ and put
\[
                         w=[[x^2,y^2],z^2]^2.          \tag{C.44.7}
\]
This word belongs to $F'\subset H$. Each of the six fillings
sends one of the three peripheral conjugacy classes to the
involution class. If $x^2$ or $y^2$ is killed, the inner
commutator in (C.44.7) is killed; if $z^2$ is killed, its
outer commutator is killed. Thus
\[
                       \varphi\sigma(w)=1
               \quad\hbox{in the full }\Gamma          \tag{C.44.8}
\]
for every one of the six maps. This is an exact free-product
identity, not a finite congruence assertion. In particular $w\in U$.

Let $q:\Gamma\to Q_{12}=\mathrm{SL}_2(\mathbb Z/12)/\{\pm I\}$
be C.23's marked quotient. The canonical signed restriction to
$F$ sends $x,y$ to $U_0^2,SU_0^2S^{-1}$, where
\[
 U_0=\begin{pmatrix}1&1\\0&1\end{pmatrix},\qquad
 S=\begin{pmatrix}0&1\\-1&0\end{pmatrix}.
\]
Modulo three the successive two commutators in (C.44.7) are
\[
 [x^2,y^2]\longmapsto\begin{pmatrix}1&1\\1&2\end{pmatrix},
 \qquad
 [[x^2,y^2],z^2]\longmapsto\begin{pmatrix}0&2\\1&0\end{pmatrix}.
                                                               \tag{C.44.9}
\]
The square of the latter is $-I$. Modulo four the two signed
images of $x^2,y^2$ are already identity, so $w$ is identity.
Consequently
\[
                         q(w)=[5I].                   \tag{C.44.10}
\]
Also $q(U)\subset q(H)=C_2^2\times Q_8$ is a two-group,
as in C.23. The executed certificate
\path{code/pbc_exact_descent_kernel_word.py}
reduces all six actual substituted words in the normal form of
$C_2*C_3$, checks the full old-marking substitution, and records
the signed matrices at three, four and twelve. Its finite matrix
checks supplement, and do not replace, (C.44.8).

\paragraph{The actual centralizer frame over a compact base.}
Put $\Gamma_2=\Gamma/O^2(U)$ and
\[
 J=\operatorname{im}\{\Gamma\to Q_{12}\times\Gamma/U\}.
\]
The map $q$ factors through $\Gamma_2$ because $q(U)$ is a
two-group. Hence there is an onto map $\Gamma_2\to J$ with
pro-two kernel $K$. The element $([5I],1)$ of $J$ exists by
(C.44.10) and centralizes the image of $r$.

Take a preliminary lift $u\in U^{(2)}$ of this element.
The order-three subgroups generated by $r$ and by $uru^{-1}$
are complements to $K$ in the preimage of the finite subgroup
$\langle r_J\rangle\simeq C_3$. They are conjugate by an
element of $K$. This is the coprime complement theorem in its
profinite form: use invariant open normal subgroups of the
pro-two kernel, apply finite coprime complement conjugacy, and
take the nonempty intersection of the compact sets of conjugators.
The two generators have the same image $r_J$, so the conjugacy
identifies the generators, not their inverses. We obtain
\[
 c\in C_{\Gamma_2}(r)\cap U^{(2)},\qquad q(c)=[5I].    \tag{C.44.11}
\]
This argument uses only the finite subgroup $C_3$ in the compact
$J$; it does not pretend that $J$ itself is finite.

In (C.44.4) take the identity automorphism with prescribed frames
\[
 b_s=b_\tau=1,\qquad
 b_r=(c\text{ at }2,\ 1\text{ at every other prime}).  \tag{C.44.12}
\]
Their common $\Gamma/U$ coordinate is one. They are genuine
centralizing frames for that identity. The closed-normal relative
lifting just proved gives an actual full positive peripheral
automorphism $\Phi$ of $\Gamma$ retaining them. Its cusp frame
lies in $R_{\rm nil}(U)$; remove it by an inner adjustment there,
so $\Phi(\tau)=\tau$ literally. All retained data persist.
Its unique pure torsion frames $A,\beta$ satisfy
\[
                     q(A)=[5I],\qquad q(\beta)=1.       \tag{C.44.13}
\]

\paragraph{Exact normalization and exact trivial descent.}
Set $f=s\Phi(s)$. Since $\Phi$ induces identity on the retained
relative quotient,
\[
        f\in R_{\rm nil}(U)\subset U'\subset F'.        \tag{C.44.14}
\]
The middle inclusion holds because every finite abelian quotient
is nilpotent. In particular the abelian normalizing exponent is
zero, not merely divisible by a finite inertia order. The full
torsion relations give the two cycle identities, and the
restriction to $F$ is the actual normalized positive gentle map
$x\mapsto x$, $y\mapsto f^{-1}yf$ in the temporary marking.

Return to the old marking
$x_{\rm old}=y^{-1}$, $y_{\rm old}=x^{-1}$.
Conjugating the restriction by $f$ fixes $x_{\rm old}$ and sends
$y_{\rm old}$ to $f y_{\rm old}f^{-1}$. Thus the normalized old
conjugator is the same abstract element $f_{\rm old}=f^{-1}$,
written in the old letters. Formula (C.44.5) is exactly the
original filling $x_{\rm old}\mapsto s$,
$y_{\rm old}\mapsto\tau$, so (C.44.6) and (C.44.14) imply
\[
                           \varphi(f_{\rm old})=1
                            \quad\hbox{in }\Gamma.     \tag{C.44.15}
\]
The actual descended automorphism fixes $s$ and sends $\tau$
to $\varphi(f_{\rm old})^{-1}\tau\varphi(f_{\rm old})$.
It is therefore the identity on the full modular group.
Its restriction is identity as well, proving $D(e)=1$ for the
full normalized gentle element $e$ just constructed.

Finally $q$ kills $R_{\rm nil}(U)$, since it factors through
$\Gamma_2$. Thus $q\Phi=q$, and the normalized cusp intertwiner
is exactly $I$ by C.12's uniqueness. The actual pure frames
(C.44.13) give the CM coordinates $([5I],1)$, so R24 E.65,
with the signed normalization checked in C.19--C.20, gives
$\Delta(e)=1$. The old-marking inner adjustment has $q(f)=1$
and changes none of these finite frames. This proves (C.44.1).

The kernel and all the character assertions now follow from
continuity and homomorphy of $D$ and $\Delta$. Since $D$ fixes
the specified Galois image, it preserves its closed normal
closure and induces an endomorphism of the arithmetic-normal
abelian quotient. The class of $e$ in that abelian quotient is
nonzero, detected by $\Delta$, and is sent to zero. This proves
the last assertion without asserting surjectivity of any descent map.
\end{proof}
